\chapter*{Part Synthesis: Applications and Modalities}
\addcontentsline{toc}{chapter}{Part Synthesis: Applications and Modalities}

\begin{whychaptermatters}
This part made the book's general workflow modality-specific. Vision sharpened invariance, language sharpened context and grounding, sequential modeling sharpened prediction-time legality, recommendation sharpened ranking, exposure, and feedback-loop reasoning, and tabular data sharpened feature encoding, leakage, calibration, and schema-drift discipline. The common lesson is that different data types still demand explicit audits before architecture choice can be trusted.
\end{whychaptermatters}

\begin{chaptersummary}
\item \textbf{Question this part taught.} What task-specific structure must be preserved, what nuisance variation should be ignored, what information was available at decision time, what did the system expose to users, what row-level and schema assumptions must hold, and what evidence makes a modality-specific claim honest under deployment conditions?
\item \textbf{Artifact chain this part added.} Vision design and stress-test card $\rightarrow$ NLP context-and-grounding card $\rightarrow$ sequential design card $\rightarrow$ recommendation audit card $\rightarrow$ tabular design card.
\item \textbf{What a student can now do.} Match invariance, context, grounding, prediction-time audits, exposure audits, and tabular leakage/monitoring discipline to the modality instead of treating pixels, tokens, sequences, ranked lists, and rows as interchangeable inputs.
\end{chaptersummary}

Three shared challenges recur across all modalities. First, benchmark success does not guarantee deployment success: shortcut learning in vision, fluent-but-false generation in language, temporal leakage in sequences, exposure bias in recommendation, and target or schema leakage in tabular systems all produce misleading offline numbers. Second, representation choices are task-specific, not universal. Third, evaluation must respect deployment constraints---camera conditions, grounding requirements, prediction-time boundaries, exposure policies, row identities, and schema contracts. The running cases (pedestrian detector, course-support assistant, wake-word detector, demand forecaster, course-resource recommender, and tabular early-warning pipeline) illustrate that the same disciplined audit applies regardless of data type.
